% !TeX root = ../papers/nc-p3.tex
\section{片側欠色の全正則窓は一つの有理曲線で覆われる}
\label{sec:klein-single-missing-windows}
\begin{theorem}\label{thm:klein-single-missing-windows}
四元群正則表現の四窓で一方の隣接切断だけに欠色があるとする。
標準14641型quadratic Koszul左右AS正則な窓は、前節のラベル変換、巡回、反転、
四位置の基底変更の下、次の一変数表示 $t\ne0$, $1+t^2\ne0$ で全て覆われる。
この表示の全点は正則で、二次PBW基底を持つ。前節の十台のうち実現する正則台は最後の一つだけである。
巡回と反転は被覆における表示操作とする。
\end{theorem}
$L=1+t^2$ と書き、各欄は $w_{j+b,j,j+e,j+e+b}$ を $j=0,1,2,3$ の順に並べた行とする。
$b=0$ は全て零である。
\[
\begin{array}{c|c|c|c}
 e & b=1 & b=2 & b=3\\\hline
 0&(-2t^2/L,-2t^2,1,1)&(t^2,1,-t^2,-1/L)&(1,-t^2,-t^2/L,1)\\
 1&(-2t^3/L,-2t,t,1/t)&(-1,1,-t^2,t^2/L)&(1,1,-t^4/L,-1)\\
 2&(0,0,0,0)&(0,0,-t^2 L,1)&(0,0,-t^2,L)\\
 3&(0,0,0,0)&(1,1,0,0)&(1,1,0,0)
\end{array}
\]
\begin{proof}[トーラス上の小行列式と有限因子]
前節の順に台を $n=0,\ldots,9$ と番号付けする。各台の非零な係数を独立なトーラス座標として
二切断の指定階数より一つ大きい全小行列式を取る。単項式を除いた二項式を先にSmith分解し、
自由トーラス上で四位置の対角基底変更をさらにSmith分解して除く。
座標の構成は定理\ref{thm:marked-klein-mixed-supports}の二段の指標と同じである。
ただし今回は $2,4$ の非自明なSmith因子を全ての根に分ける。
残りの小行列式を代入して非零変数の積を局所化すると、各台で有限因子の選択が一通りだけ残る。
他の根では局所化したイデアルが単位になる。この検査は $\mathbb Q(\sqrt{-1})$ 上で行える。
残る変数を $a,b,c,d,e,f,g$ の必要な個数だけ左から取ると、以下は残る式による
稠密な有理表示である。右辺は上から順に代入する。$H$ は始点0、長さ4、四元群総次数0の次元である。
\[
\begin{array}{c|l|l|r}
 n&\text{自由変数}&\text{他変数}& H\\\hline
0&a,d&b=ad^2+d,\ c=ad/(a-b)&12\\
1&a& &12\\
2&a,b&c=1+1/b&12\\
3&b,c&a=1/b-1/c&12\\
4&c,d,f&a=f/(d(f-1)),\ e=1-ad^2/(cf),\ b=1+ce/a&12\\
5&a,b&c=1+ab&13\\
6&b,c,e&a=e(bc-1)/(b-1),\ d=(1-e)/(c-e)&12\\
7&b,d,e,f&c=def-d+1,\ a=(c-f)/(e-1)&12\\
8&b,c,d,f&e=bf+cf-bcf^2,\ g=(bd^2-e)/(bd^2-1),\ a=bd(c-g)&12\\
9&d&a=-d^2,\ b=1/(a-1),\ c=-2a^3b/d^3&11
\end{array}
\]

分母によって別成分を落としていないことを示す。0--5と9では各式の分母の消失は
元の小行列式と座標の非零性に反する。例えば0で $a=b$ なら $ad=0$、
4で $f=1$ なら $f=0$、9で $a=1$ なら $b(a-1)=1$ に反する。
6の方程式は $(b-1)a=e(bc-1),(c-e)d=1-e$ である。
$b=1$ なら $c=1$、$c=e$ なら $c=e=1$ であり、どちらの境界も次元2以下である。
五変数の二式の各成分は次元3以上なので新しい成分はない。
7で $e=1$ なら $c=f,(d-1)(f-1)=0$ で次元3以下、六変数の二式は次元4以上である。
8では $a,e$ を消去した後の五変数の一式で $bd^2=1$ なら $e=1$ と
$(1-bf)(1-cf)=0$ が必要になり次元3以下である。一般開集合の次元は4であり、これも別成分を持たない。
各行のグラフの残る非零開集合は既約であるから、表の計算は全実現点を含む閉包上の一般点計算になっている。

四次障害を再現するには、総次数零の64語を列に取り、三位置の二次関係の全挿入を行とし、
表の自由変数の有理関数体で余核次元を求める。各色ブロックの実際の階数が指定値に等しいところでは
切断像はベクトル束であり、余核次元は特殊化で下がらない。
定理\ref{thm:klein-missing-color-coverage}の漸化式では、始点0の各非自明指標の長さ4での値が3なので
標準値は $(35+3\cdot3)/4=11$ である。これで初めの九台は例外値を含め排除される。
9の有限因子は $2,4$ の二つであり、残った根で $d=t$ とすれば上の係数表そのものになる。
\path{computations/verify_klein_single_missing_windows.py} は全小行列式、全ての根、
自由変数体のグラフと四次階数を検算する。
\end{proof}

\begin{proof}[二つの二文字窓の線型旗]
位置0,1では第一組 $A=\{0,1\}$、第二組 $B=\{2,3\}$ とし、位置2,3では二組を逆にする。
各辺に第一組内部と第二組内部を結ぶ関係が一つずつあり、ともに可逆な二行二列の双線型形式である。
残る四関係は $A_p B_{p+1}+B_p A_{p+1}$ に入り、$B_p A_{p+1}$ への射影が可逆である。
実際、係数表の各ブロックで辞書順の先頭独立列を取ると、
この混合四行列の行列式は四位置で順に
\[
           L^2,\qquad 4t^6/L^2,\qquad -t^4 L^2,\qquad -16t^8/L^2 .
\]
二つの内部形式の行列式も各位置で非零定数と $t,L$ の冪だけからなる。
これは上記検算の最後の四つの射影で全値について確認できる。

内部形式により、第一組の一文字の射影的評価から次位置の関係を消す相手の評価への
射影直線の同型が定まる。四位置を一巡する合成は射影直線の自己同型なので固定点を持つ。
その点から四位置を一周して、各位置の第一組の最大文字に対応する双対基底ベクトルを選ぶ。
内部関係で最大文字同士の係数が消え、次の列は非零なので先頭対は二文字の下降対になる。
第二組にも同じ選択を独立に行う。各位置で第一組を第二組より小さく置くと、
混合四関係は全ての $BA$ を $AB$ に変える。これで全ての位置で通常の六下降対が得られた。
二十四項全非零をここでは必要とせず、係数表から一文字縮約が全位置で単射である。
三次元数の上下界20による重なり解消、Koszul双対の完全対合は定理\ref{thm:unmarked-klein-pbw}と同じである。
従って全点で左右AS正則である。
\end{proof}
